% ID: FIS-DIN-ATT-001
% Titolo: Cassa spinta su un pavimento ruvido
% Disciplina: Fisica
% Area: Dinamica
% Argomento: Attrito
% Sottoargomento: Attrito dinamico e secondo principio
% Classe: Terza liceo scientifico
% Difficolta: 3
% Tipo: problema_numerico
% Risultato: F_attrito = 49 N; a = 2,05 m/s^2; v = 8,2 m/s; s = 16,4 m
% Tempo_stimato: 10 min
% Competenze: bilancio delle forze, applicazione del secondo principio, cinematica
% Prerequisiti: forza di attrito, forza peso, moto uniformemente accelerato
% Tag: fisica, dinamica, attrito, secondo_principio, moto_accelerato
% Fonte: esempio_originale
% Autore: OpenAI Codex
% Licenza: CC-BY-SA-4.0

\begin{esercizio}
Una cassa di massa \(20\,\mathrm{kg}\) e spinta su un pavimento orizzontale
da una forza costante di modulo \(90\,\mathrm{N}\), parallela al pavimento.
Il coefficiente di attrito dinamico tra cassa e pavimento e
\(\mu_d=0{,}25\).

Supponendo che la cassa parta da ferma, determina l'accelerazione e poi la
velocita e lo spazio percorso dopo \(4{,}0\,\mathrm{s}\). Usa
\(g=9{,}8\,\mathrm{m/s^2}\).
\end{esercizio}

\begin{soluzione}
Sul piano orizzontale la reazione normale vale:
\[
N=mg=20\cdot 9{,}8=196\,\mathrm{N}.
\]

La forza di attrito dinamico ha modulo:
\[
F_\text{attrito}=\mu_d N=0{,}25\cdot 196=49\,\mathrm{N}.
\]

La forza risultante lungo la direzione del moto e:
\[
F_\text{tot}=90-49=41\,\mathrm{N}.
\]

Dal secondo principio della dinamica:
\[
a=\frac{F_\text{tot}}{m}=\frac{41}{20}=2{,}05\,\mathrm{m/s^2}.
\]

Partendo da ferma, dopo \(4{,}0\,\mathrm{s}\) la velocita e:
\[
v=at=2{,}05\cdot 4{,}0=8{,}2\,\mathrm{m/s}.
\]

Lo spazio percorso e:
\[
s=\frac{1}{2}at^2
=\frac{1}{2}\cdot 2{,}05\cdot (4{,}0)^2
=16{,}4\,\mathrm{m}.
\]
\end{soluzione}
